% =====================================================================
%  Chapter 8 — Conclusion and Future Work  (~900 words)
%  DRAFT generated 2026-07-18.
% =====================================================================
\chapter{Conclusion and Future Work}
\label{ch:conclusion}

\section{Summary of Contributions}
\label{sec:conc-summary}

This dissertation asked when an offline, CPU-only medical question-answering
system should retrieve, and answered it with a training-free three-gate
ensemble evaluated end-to-end on a quantised 3B model over a 426K-chunk
medical corpus. The contributions, with their evidence:

\textbf{A working training-free gate} (C1). Token entropy, logit margin and
a two-draft hallucination probe, majority-voted, with hardware-aware
degradation --- no fine-tuning, no labelled data, no second model
(Chapter~\ref{ch:method}).

\textbf{Calibration non-transfer} (C2). Literature thresholds are physically
unreachable on the quantised model (observed entropy 0.27--1.59 nats against
a 2.5-nat trigger); offline replay recalibration finds usable operating
points in seconds. Gate thresholds are model-specific and must be treated as
such (Section~\ref{sec:eval-calibration}).

\textbf{Retrieval as distraction} (C3). Retrieval broke 46 previously
correct answers and fixed 13 --- harmful $3.5{:}1$ --- with chunks equally
on-topic but $2.4\times$ less answer-bearing where it harmed
(Section~\ref{sec:eval-distraction}).

\textbf{The protective skip} (C4). On those poisoned questions the gate's
skip scored 19/19; its retrieve, 3/27. The gate's value on this evidence is
protection at the boundary of parametric knowledge, not merely cost saving
(Section~\ref{sec:eval-protective}).

\textbf{The safety envelope} (C5). No policy clears any stipulated clinical
bar; the best, closed-book at \accPthreeMcq{}, is \gapToLowBar{}pp short of
the lowest. The measured distance to clinical acceptability for this
deployment class is the honest deliverable
(Section~\ref{sec:eval-envelope}).

\textbf{A reproducible harness} (C6). Log-generated numbers with a
drift-guard test, 180+ model-free unit tests, paired-run proofs, exact
determinism (Chapter~\ref{ch:impl}).

Every claim is held at the strength its evidence supports: the retrieval
halving (\retPfiveMcq{} versus \retPoneMcq{}) is exact; the accuracy gain
(\diffPfivePone{}pp, 95\% CI \diffPfivePoneCI{}) is consistent with an
improvement but not established at $n=\nQuestions{}$; the envelope is empty.
If one sentence survives this dissertation, it should be this: retrieval,
the field's default remedy for hallucination, made this model \emph{worse}
on 46 of 200 questions it already knew --- and the gate's contribution is
knowing when not to use it.

\section{Future Work}
\label{sec:conc-future}

Ordered by ratio of value to effort, each specific enough to pick up:

\begin{enumerate}
  \item \textbf{Cross-model scaling study.} All prerequisites exist and are
    verified: a chat-format abstraction proven byte-identical for existing
    logs, Qwen-2.5 7B/14B configurations at matched quantisation, a
    four-check smoke test gating long runs, and a resumable GPU notebook.
    Running P3+P5 (MCQ and open-ended) on one larger model answers the two
    questions single-model evidence cannot: does calibration non-transfer
    recur (expected: yes, with a different operating point), and does
    selective-beats-always survive scale? It also tests whether the
    closed-book ceiling rises toward the safety bars.
  \item \textbf{Content-aware entropy.} Section~\ref{sec:eval-attribution}
    showed raw token entropy fires on phrasing, not clinical content.
    Restricting the mean to content tokens, or weighting by attention to the
    question's clinical terms, directly attacks the method's clearest
    measured weakness --- and can be prototyped offline against the logged
    per-token arrays before any live run.
  \item \textbf{Retrieval-aware (post-hoc) gating.} The distraction analysis
    showed the harmful case is topically-relevant, non-answer-bearing
    chunks. The current gate must commit \emph{before} seeing any evidence;
    a second-stage gate asking ``having seen the chunks, are they
    decisive?'' could reject a retrieval after the fact and reclaim part of
    the 3/27 failure region.
  \item \textbf{Gate cost optimisation.} The entropy and margin signals
    derive from the same draft but are currently generated twice; merging
    them cuts five model calls per query to four ($\sim$20\% latency), and a
    probe-free logit-only configuration would remove the two probe drafts
    that dominate cost --- at a measured accuracy trade-off to be
    quantified. Section~\ref{sec:eval-cost} shows latency, not retrieval
    rate, is the binding cost on CPU.
  \item \textbf{Statistical power.} The full 7,663-question MIRAGE
    benchmark, multiple seeds, and a purposely balanced high-risk stratum
    --- the current $n=6$ is the single weakest point of the safety
    analysis --- would convert suggestive intervals into decidable ones.
  \item \textbf{Deployment surface.} A human-in-the-loop interface exposing
    the gate decision, its signals and per-token uncertainty to a clinician
    reviewer; and physical benchmarking of the low/high hardware tiers on
    real 4\,GB and 16\,GB devices, closing the gap between the implemented
    mechanism and the measured claim.
\end{enumerate}

The through-line of all six is the same as the dissertation's: in
safety-critical, resource-constrained settings, the interesting question is
not how to make a model retrieve better, but how to make a system know when
its model already knows.
